\chapter{The bus}
\label{ch:buss}

{\large\itshape Two wires, many nodes, and an asymmetry that decides everything.\par}

\begin{chapterbox}{In this chapter}
\begin{itemize}
  \item Why point-to-point stops working as a system grows.
  \item The shared bus, and what it costs.
  \item Dominant and recessive bits, and the logical AND the bus performs.
  \item Differential signalling and termination, as far as someone building a node needs them.
\end{itemize}
\end{chapterbox}

\section{The problem}

An embedded system with five units that all need to talk to each other needs ten connections if
every pair is wired together directly. With ten units it is forty-five. The formula is $n(n-1)/2$,
and it is relentless: the wiring grows quadratically while the benefit grows linearly.

In a car that is not a theoretical objection. Engine management, ABS, instrument cluster,
gearbox, airbags, climate control --- and every connection is copper that weighs, costs, takes up
space and can break.

\textbf{The solution is to stop wiring units to each other, and wire them in instead.} A single
bus, shared by every node:

\bookfigure{bus_topology}{Point-to-point versus a shared bus. Five nodes need ten connections in
the case on the left and one bus in the case on the right.}{fig:topology}

What was $n(n-1)/2$ connections becomes $n$ attachments to one common line. The price is that the
bus is now a shared resource, and shared resources have to be allocated.

\section{What a shared bus costs}

Three problems appear the moment two nodes share a wire, and the whole design of CAN is the answer
to them:

\begin{booktable}{@{}lL@{}}
\thead{Problem} & \thead{CAN's answer} \\ \midrule
Everyone hears everything, so a receiver must know what concerns it & The identifier describes
  the \emph{content}, and every node filters for itself (\kapref{ch:frameformat}) \\
Two nodes can start sending at the same time & Non-destructive arbitration (\kapref{ch:arbitrering}) \\
A disturbance on the wire hits everyone & Bit stuffing and CRC-15 (\kapref{ch:stuffing}), plus
  error handling (\kapref{ch:fel}) \\
\end{booktable}

Notice what is \emph{not} in the table: addressing. CAN has no destination address. That is not a
simplification but a design choice, and \kapref{ch:frameformat} goes through why.

\section{Dominant and recessive}

Here is the one idea you need to have entirely clear for the rest of CAN to make sense.

A CAN bus has two states, and they are \textbf{not} symmetric:

\begin{narrowtable}{@{}lll@{}}
\thead{State} & \thead{Logic level} & \thead{Wins?} \\ \midrule
\dom{} & 0 & Yes \\
\rec{} & 1 & No \\
\end{narrowtable}

The rule is simple and has consequences all through the book:

\begin{aside}
\textbf{The bus is dominant if \emph{at least one} node drives it dominant. It is recessive only if
\emph{every} node lets go of it.}
\end{aside}

Electrically it is an open drain with a pull-up resistor: a node can only \emph{pull} the line
down, never drive it up. If everyone lets go, it is pulled up on its own.

Logically, then, the bus is an AND gate over the outputs of every node:

\bookfigure{wired_and}{The bus as a wired AND. Three nodes, and the bus level as a function of what
they drive: a single dominant node is enough to make the bus dominant.}{fig:wiredand}

\subsection{Why the asymmetry is needed}

Suppose instead a symmetric bus, where a one and a zero that meet give an undefined level in
between. Then a collision is \emph{destruction}: both frames become unusable, and both nodes have to
detect it, wait for different lengths of time and try again. That is how Ethernet did it, and it is
why Ethernet needed random backoff.

With the asymmetry, a collision becomes a \emph{comparison} instead. The node that sends recessive
and reads back dominant immediately knows two things: that someone else is sending too, and that the
other one has priority. It withdraws, silently, in the middle of the frame --- and the winning node's
frame continues without losing a single bit interval.

This is called \textbf{non-destructive arbitration}, and it is CAN's most elegant property.
\Kapref{ch:arbitrering} works through it bit by bit.

\subsection{Reading back what you send}

A precondition that is easy to overlook: \textbf{every transmitting node reads the bus continuously
while it sends.} That is how it detects that it has lost arbitration, and it is also how it detects
that the bus is not doing what it was told.

The simplified controller does this during the arbitration field. A complete CAN controller does it
during the \emph{whole} frame, and reports a \emph{bit error} as soon as what it reads back differs
from what it sent. The difference is described in \kapref{ch:kursen}.

\section{The wires}

CAN is \textbf{differential}: the signal is carried as a voltage difference between two conductors,
\code{CAN_H} and \code{CAN_L}, twisted around each other.

\begin{narrowtable}{@{}lll@{}}
\thead{State} & \thead{\code{CAN_H}} & \thead{\code{CAN_L}} \\ \midrule
\rec{} & 2.5 V & 2.5 V \\
\dom{} & 3.5 V & 1.5 V \\
\end{narrowtable}

In the recessive state nobody drives, and both conductors rest at the same level --- the difference
is zero. In the dominant state the transmitting node drives them apart, to about two volts.

The point of differential signalling is noise immunity. A disturbance that hits the cable hits both
conductors about equally, and since the receiver measures only the \emph{difference} it takes no
notice of it. The twisting is what makes the disturbance really hit both equally.

\subsection{Termination}

A CAN bus is terminated with \textbf{120~\textohm{} at each end}, not in the middle, and never more
than two resistors no matter how many nodes are on the bus.

The resistors match the cable's characteristic impedance, so that a signal reaching the end is
absorbed instead of reflecting back. A badly terminated bus often works anyway on a short cable at
low speed, and stops working just when you have started to trust it --- which makes it one of the
nastier faults.

\begin{aside}
\note[At the bench] A quick way to check the termination: measure the resistance across
\code{CAN_H} and \code{CAN_L} with the bus unpowered. Two 120~\textohm{} in parallel should give
about 60~\textohm. If you get 120, a resistor is missing; if you get 40, there are three.
\end{aside}

\section{Speed and length}

CAN's bit rate and the length of the bus are linked, and the reason is arbitration: during the
arbitration field, a bit must have time to reach the farthest node on the bus \emph{and} come back
before the bit ends, or the nodes cannot agree on who won.

\begin{narrowtable}{@{}ll@{}}
\thead{Bit rate} & \thead{Approximate maximum length} \\ \midrule
1 Mbit/s & 40 m \\
500 kbit/s & 100 m \\
125 kbit/s & 500 m \\
\end{narrowtable}

The simplified controller runs at \textbf{1 Mbit/s}, and the bus is a cable on the lab bench. The
limit is therefore never anywhere near becoming a problem --- but the relationship is worth knowing,
because it explains why vehicle systems often run slower than they could.

\section*{Exercises}

\exercise{Wiring}{Calculation}
\label{ex:1:1}
\begin{enumerate}[a)]
  \item How many connections are needed to wire eight units together point-to-point?
  \item How many attachments are needed if they share a bus instead?
  \item At what number of units does point-to-point need twice as many connections as the bus has
    attachments?
\end{enumerate}

\exercise{The bus level}{Reasoning}
\label{ex:1:2}
Four nodes are on a bus. During a certain bit interval, node A drives recessive, node B dominant,
node C recessive, and node D does not drive at all.
\begin{enumerate}[a)]
  \item What level is the bus at?
  \item Which of the nodes read back something other than what they sent?
  \item What conclusion do they draw from that?
\end{enumerate}

\exercise{Termination}{Reasoning}
\label{ex:1:3}
A group measures 40~\textohm{} across an unpowered bus.
\begin{enumerate}[a)]
  \item What is wrong?
  \item The bus works anyway on the lab bench. Why, and when would it stop doing so?
\end{enumerate}

\exercise{The asymmetry}{Reasoning}
\label{ex:1:4}
Imagine a hypothetical bus where recessive won over dominant instead --- that is, where the bus is
recessive if at least one node drives it recessive.
\begin{enumerate}[a)]
  \item Would arbitration still work?
  \item Which identifier value would then have the highest priority?
  \item Why did CAN choose the order it did? Think about what the bus does when nobody is sending.
\end{enumerate}
